\chapter{2024年新高考I卷}

\section{单选题}

\begin{problem}
已知 \(A=\{x\mid -\sqrt[3]{5}<x<\sqrt[3]{5}\}\)，\(B=\{-3,-1,0,2,3\}\)，则 \(A\cap B=\)（\quad）

\choices
    {\(\{-1,0\}\)}
    {\(\{2,3\}\)}
    {\(\{-3,-1,0\}\)}
    {\(\{-1,0,2\}\)}
\end{problem}

\begin{answer}
A
\end{answer}

\begin{solution}
因为 \(1<\sqrt[3]{5}<2\)，所以
\[
-\sqrt[3]{5}<-1<0<\sqrt[3]{5}
\]
而 \(-3\)，\(2\)，\(3\) 都不在区间 \(\left(-\sqrt[3]{5},\sqrt[3]{5}\right)\) 内，因此
\(A\cap B=\{-1,0\}\).
故选 A.
\end{solution}

\begin{problem}
若 \(\dfrac{z}{z-1}=1+\i\)，则 \(z=\)（\quad）

\choices
    {\(-1-\i\)}
    {\(-1+\i\)}
    {\(1-\i\)}
    {\(1+\i\)}
\end{problem}

\begin{answer}
C
\end{answer}

\begin{solution}
由
\(\frac{z}{z-1}=1+\i\)
得
\(z=(1+\i)(z-1)\).
所以
\(\i z=-(1+\i)\)，
从而
\(z=\frac{-(1+\i)}{\i}=1-\i\).
故选 C.
\end{solution}

\begin{problem}
已知向量 \(\bs a=(0,1)\)，\(\bs b=(2,x)\)，若 \(\bs b\perp (\bs b-4\bs a)\)，则 \(x=\)（\quad）

\choices
    {\(-2\)}
    {\(-1\)}
    {\(1\)}
    {\(2\)}
\end{problem}

\begin{answer}
D
\end{answer}

\begin{solution}
由 \(\bs b\perp (\bs b-4\bs a)\)，得
\(\bs b\cdot (\bs b-4\bs a)=0\).
即
\((2,x)\cdot\left(2,x-4\right)=4+x(x-4)=0\).
故 \(x^2-4x+4=0\)，所以 \(x=2\)（题中仅有一个取值）. 故选 D.
\end{solution}

\begin{problem}
已知 \(\cos(\alpha+\beta)=m\)，\(\tan\alpha\tan\beta=2\)，则 \(\cos(\alpha-\beta)=\)（\quad）

\choices
    {\(-3m\)}
    {\(-\frac{m}{3}\)}
    {\(\frac{m}{3}\)}
    {\(3m\)}
\end{problem}

\begin{answer}
A
\end{answer}

\begin{solution}
由 \(\cos(\alpha+\beta)=m\) 得
\(\cos\alpha\cos\beta-\sin\alpha\sin\beta=m\).
又 \(\tan\alpha\tan\beta=2\)，故
\(\sin\alpha\sin\beta=2\cos\alpha\cos\beta\).
于是
\(\cos\alpha\cos\beta-2\cos\alpha\cos\beta=-m\)
即 \(\cos\alpha\cos\beta=-m\). 从而
\[
\cos(\alpha-\beta)=\cos\alpha\cos\beta+\sin\alpha\sin\beta=(-m)+2(-m)=-3m
\]
故选 A.
\end{solution}

\begin{problem}
已知圆柱和圆锥的底面半径相等，侧面积相等，且它们的高均为 \(\sqrt3\)，则圆锥的体积为（\quad）

\choices
    {\(2\sqrt3\pi\)}
    {\(3\sqrt3\pi\)}
    {\(6\sqrt3\pi\)}
    {\(9\sqrt3\pi\)}
\end{problem}

\begin{answer}
B
\end{answer}

\begin{solution}
设圆柱和圆锥的底面半径都为 \(r\). 圆锥的母线长为
\(l=\sqrt{r^2+(\sqrt3)^2}=\sqrt{r^2+3}\).
侧面积相等给出
\(2\pi r\cdot \sqrt3=\pi rl\)
即
\(\sqrt{r^2+3}=2\sqrt3\)，
故
\(r^2=9\).
所以 \(r=3\). 圆锥体积
\[
V=\frac13\pi r^2\cdot \sqrt3=\frac13\pi\cdot 9\cdot \sqrt3=3\sqrt3\pi
\]
故选 B.
\end{solution}

\begin{problem}
已知函数
\(f(x)=-x^2-2ax-a\)（\(x<0\)），\(f(x)=\e^x+\ln(x+1)\)（\(x\ge 0\)）
在 \(\R\) 上单调递增，则 \(a\) 取值范围为（\quad）

\choices
    {\((-\infty,0]\)}
    {\([-1,0]\)}
    {\([-1,1]\)}
    {\([0,+\infty)\)}
\end{problem}

\begin{answer}
B
\end{answer}

\begin{solution}
当 \(x\ge 0\) 时，
\(f(x)=\e^x+\ln(x+1)\)，
故
\(f'(x)=\e^x+\frac1{x+1}>0\)，
于是右侧部分严格递增.

对 \(x<0\)，有
\(f'(x)=-2x-2a\).
要使该部分单调递增需 \(f'(x)\ge 0\)，即 \(-2x-2a\ge 0\)，对 \(x<0\) 最严格在 \(x\to0^-\) 时取到，故 \(-2a\ge 0\)，即 \(a\le0\).

此外，单调性要求左侧极限不超过右侧端点值，即
\[
\lim\limits_{x\to0^-}f(x)=-a\le f(0)=\e+\ln1=1
\]
故 \(-a\le1\)，即 \(a\ge-1\). 故 \(a\in[-1,0]\).
\end{solution}

\begin{problem}
当 \(x\in[0,2\pi]\) 时，曲线 \(y=\sin x\) 与 \(y=2\sin\!\left(3x-\frac{\pi}{6}\right)\) 的交点个数为（\quad）

\choices
    {3}
    {4}
    {6}
    {8}
\end{problem}

\begin{answer}
C
\end{answer}

\begin{solution}
设
\(h(x)=\sin x-2\sin\left(3x-\frac{\pi}{6}\right)\).
题目所求即为方程 \(h(x)=0\) 在 \([0,2\pi]\) 上解的个数.

因为
\(h(x+\pi)=\sin(x+\pi)-2\sin\left(3x+\frac{17\pi}{6}\right)=-h(x)\)，
所以 \(h(x)=0\) 的解关于平移 \(\pi\) 成对出现，只需考察区间 \([0,\pi]\).

当 \(x\in\left[\frac{\pi}{6},\frac{\pi}{3}\right]\) 时，\(\sin x\le \frac{\sqrt3}{2}\)，且 \(3x-\frac{\pi}{6}\in\left[\frac{\pi}{3},\frac{5\pi}{6}\right]\)，从而 \(\sin\left(3x-\frac{\pi}{6}\right)\ge \frac12\). 于是
\(h(x)\le \frac{\sqrt3}{2}-1<0\).
当 \(x\in\left[\frac{\pi}{2},\frac{2\pi}{3}\right]\) 时，\(\sin x\ge \frac{\sqrt3}{2}\)，且 \(3x-\frac{\pi}{6}\in\left[\frac{4\pi}{3},\frac{11\pi}{6}\right]\)，从而 \(\sin\left(3x-\frac{\pi}{6}\right)\le -\frac12\). 于是
\(h(x)\ge \frac{\sqrt3}{2}+1>0\).
当 \(x\in\left[\frac{5\pi}{6},\pi\right]\) 时，\(\sin x\le \frac12\)，且 \(3x-\frac{\pi}{6}\in\left[\frac{7\pi}{3},\frac{17\pi}{6}\right]\)，其正弦值不小于 \(\frac12\)，于是
\(h(x)\le \frac12-1<0\).

\solutionfigure{\bitmapfigure[max width=6.8cm]{img/2024/new_gaokao_paper_1/q7_fig1.png}}

再看三个剩余区间：
\(h(0)=1>0\)，\(h\left(\frac{\pi}{6}\right)=\frac12-\sqrt3<0\)，
\(h\left(\frac{\pi}{3}\right)=\frac{\sqrt3}{2}-1<0\)，\(h\left(\frac{\pi}{2}\right)=1+\sqrt3>0\)，
\(h\left(\frac{2\pi}{3}\right)=1+\frac{\sqrt3}{2}>0\)，\(h\left(\frac{5\pi}{6}\right)=\frac12-\sqrt3<0\).
由介值定理，方程 \(h(x)=0\) 在
\(\left(0,\frac{\pi}{6}\right)\)，\(\left(\frac{\pi}{3},\frac{\pi}{2}\right)\)，\(\left(\frac{2\pi}{3},\frac{5\pi}{6}\right)\)
内各有一个解，而在其余部分没有解，所以在 \([0,\pi]\) 上共有 \(3\) 个解.

再由 \(h(x+\pi)=-h(x)\)，得到 \([\pi,2\pi]\) 上对应地也有 \(3\) 个解. 因此两个曲线在 \([0,2\pi]\) 上共有 \(6\) 个交点. 故选 C.
\end{solution}

\begin{problem}
已知函数 \(f(x)\) 的定义域为 \(\R\)，且
\(f(x)>f(x-1)+f(x-2)\)，
当 \(x<3\) 时 \(f(x)=x\). 则下列结论中一定正确的是（\quad）

\choices
    {\(f(10)>100\)}
    {\(f(20)>1000\)}
    {\(f(10)<1000\)}
    {\(f(20)<10000\)}
\end{problem}

\begin{answer}
B
\end{answer}

\begin{solution}
由题设得 \(f(1)=1\)，\(f(2)=2\). 由递推式 \(f(x)>f(x-1)+f(x-2)\) 可得 \(f(3)>3\)，\(f(4)>5\)，\(f(5)>8\)，\(f(6)>13\)，\(\dots\)，逐项递增且均大于对应斐波那契数，故
\(f(10)>89\)，\(f(20)>1597>1000\).
故 \(f(20)>1000\) 正确，其他不一定. 故选 B.
\end{solution}
\section{多选题}

\begin{problem}
为了解推动出口后的亩收入（单位：万元）情况，从该种植区抽取样本，得到推动出口后亩收入样本均值 \(x=2.1\)，样本方差 \(s^2=0.01\). 已知该种植区以往亩收入 \(X\sim N(1.8,0.1)\)，假设推动出口后的亩收入 \(Y\sim N(x,s)\)，则（\quad）（若 \(Z\sim N(u,\sigma)\)，则 \(P(Z<u+\sigma)\approx0.8413\)）

\choices
    {\(P(X>2)>0.2\)}
    {\(P(X>2)<0.5\)}
    {\(P(Y>2)>0.5\)}
    {\(P(Y>2)<0.8\)}
\end{problem}

\begin{answer}
BC
\end{answer}

\begin{solution}
由 \(x=2.1\)，\(s^2=0.01\)，有 \(Y\sim N(2.1,0.1)\). 故
\[
P(Y>2)=P(Y>2.1-0.1)=P(Y<2.1+0.1)\approx0.8413>0.5
\]
所以 \(C\) 正确，\(D\) 错误.

又 \(X\sim N(1.8,0.1)\)，故
\[
P(X>2)=P\left(X>1.8+2\times0.1\right)<P\left(X>1.8+0.1\right)\approx1-0.8413=0.1587<0.2
\]
所以 B 正确，A 错误. 故选 BC.
\end{solution}

\begin{problem}
设函数
\(f(x)=(x-1)^2(x-4)\)
，则（\quad）

\choices
    {\(x=3\) 是 \(f(x)\) 的极小值点}
    {当 \(0<x<1\) 时，\(f(x)<f(x^2)\)}
    {当 \(1<x<2\) 时，\(-4<f(2x-1)<0\)}
    {当 \(-1<x<0\) 时，\(f(2-x)>f(x)\)}
\end{problem}

\begin{answer}
ACD
\end{answer}

\begin{solution}
设 \(x\in\R\)，
\(f'(x)=3(x-1)(x-3)\).
当 \(x\in(-\infty,1)\cup(3,+\infty)\) 时 \(f'(x)>0\)，当 \(x\in(1,3)\) 时 \(f'(x)<0\)，故 \(x=1\) 是局部极大，\(x=3\) 为局部极小；在 \(0<x<1\) 时 \(f\) 增，故 \(f(x)<f(x^2)\) 的说法方向相反，\(B\) 错误.
当 \(1<x<2\) 时，\(2x-1\in(1,3)\) 且 \(f\) 在 \((1,3)\) 上递减，所以 \(f(1)>f(2x-1)>f(3)\)，即 \(-4<f(2x-1)<0\)，故 \(C\) 正确.
当 \(-1<x<0\)，
\(f(2-x)-f(x)=2(1-x)^3\).
故 \(f(2-x)>f(x)\)，所以 D 正确. 故选 ACD.
\end{solution}

\begin{problem}
造型可做成美丽的丝带，将其看作图中曲线 \(C\) 的一部分. 已知 \(C\) 过坐标原点 \(O\)，且 \(C\) 上点满足横坐标大于 \(-2\)，到点 \(F(2,0)\) 的距离与到定直线 \(x=a\)（\(a<0\)） 的距离之积为 \(4\)，则（\quad）

\bitmapfigure[max width=6.0cm]{img/2024/new_gaokao_paper_1/q11_fig1.png}

\choices
    {\(a=-2\)}
    {点 \((2\sqrt2,0)\) 在 \(C\) 上}
    {在第一象限 \(C\) 的点的纵坐标最大值为 \(1\)}
    {若点 \((x_0,y_0)\) 在 \(C\) 上，则 \(y_0\le \dfrac{4}{x_0+2}\)}
\end{problem}

\begin{answer}
ABD
\end{answer}

\begin{solution}
设曲线上点为 \(P(x,y)\)，则 \(x>-2\) 且
\(\sqrt{(x-2)^2+y^2}\cdot |x-a|=4\).
又 \(C\) 过 \(O\)，代入得 \(2\cdot|0-a|=4\)，故 \(a=-2\).
因此曲线方程可写为
\(\sqrt{(x-2)^2+y^2}\,(x+2)=4\).
两边平方得
\(\left((x-2)^2+y^2\right)(x+2)^2=16\).
当 \(x=2\sqrt2\)，\(y=0\) 时，
\(\sqrt{(2\sqrt2-2)^2+0^2}\,(2\sqrt2+2)=2(\sqrt2-1)\cdot 2(\sqrt2+1)=4\)，
故点 \((2\sqrt2,0)\) 在 \(C\) 上，\(B\) 正确.
又当 \(x=\frac32\) 时，
\[
y^2=\frac{16}{(x+2)^2}-(x-2)^2=\frac{64}{49}-\frac14=\frac{207}{196}>1
\]
故第一象限上 \(y\) 的最大值大于 \(1\)，所以 \(C\) 错.
再由方程可得
\(y_0^2=\frac{16}{(x_0+2)^2}-(x_0-2)^2\).
于是
\(y_0^2\le \frac{16}{(x_0+2)^2}\).
因为 \(x_0+2>0\)，所以
\(y_0\le \frac{4}{x_0+2}\).
故 \(D\) 正确.
\end{solution}
\section{填空题}

\begin{problem}
设双曲线
\(\frac{x^2}{a^2}-\frac{y^2}{b^2}=1\ (a>0,b>0)\)
的左右焦点分别为 \(F_1\)，\(F_2\). 过 \(F_2\) 作与 \(y\) 轴平行的直线交该双曲线于 \(A\)，\(B\) 两点. 若
\(|AF_1|=13\)，\(|AB|=10\)，
则该双曲线的离心率为（\quad）.

\bitmapfigure[max width=6.0cm]{img/2024/new_gaokao_paper_1/q12_fig1.png}
\end{problem}

\begin{answer}
\(\dfrac{3}{2}\)
\end{answer}

\begin{solution}
设两交点横坐标均为 \(c\). 由方程可得
\(y_A=\frac{b}{a}\sqrt{c^2-a^2}\)，\(y_B=-\frac{b}{a}\sqrt{c^2-a^2}\).
故 \(|AB|=2\frac{b}{a}\sqrt{c^2-a^2}=10\)，即 \(AB=10\Rightarrow \frac{b}{a}\sqrt{c^2-a^2}=5\)，并且
\(|AF_2|=5\).
又 \(F_2\) 到 \(F_1\) 的距离为 \(2a\)，所以
\(|AF_1|=|AF_2|+2a=13\Rightarrow a=4\).
代入 \(|AB|=10\) 可得 \(b^2=20\). 故离心率
\(e=\sqrt{1+\frac{b^2}{a^2}}=\sqrt{1+\frac{20}{16}}=\frac32\).
\end{solution}

\begin{problem}
若曲线 \(y=\e^x+x\) 与曲线 \(y=\ln(x+1)+a\) 在点 \((0,1)\) 处有切线重合，则 \(a=\)\fillinblank{}.
\end{problem}

\begin{answer}
\(\ln2\)
\end{answer}

\begin{solution}
\(y=\e^x+x\) 在 \((0,1)\) 的切线方程是
\(y=2x+1\).
设 \(y=\ln(x+1)+a\) 在 \(x_0\) 处切于该点且切线与上述相同. 因为该切线斜率是 2，故
\(\frac{1}{x_0+1}=2\Rightarrow x_0=-\frac12\).
因此另一切点坐标为 \(\left(-\frac12,a+\ln\frac12\right)\)，其切线
\(y=2x+1+a-\ln2\).
与 \(y=2x+1\) 重合得 \(a=\ln2\).
\end{solution}

\begin{problem}
甲，乙两人各有四张卡片，甲为 \(\{1,3,5,7\}\)，乙为 \(\{2,4,6,8\}\). 每轮各出一张牌比较大小，先大者得 \(1\) 分，先小者得 \(0\) 分，弃置该轮牌. 四轮后甲总得分不少于 \(2\) 的概率是\fillinblank{}.
\end{problem}

\begin{answer}
\(\frac12\)
\end{answer}

\begin{solution}
把甲四张卡片的出牌顺序固定，乙的四张卡片按随机顺序依次对应，则共有 \(A_4^4=4!\) 种等可能结果.
设甲四轮得分分别为 \(X_1\)，\(X_2\)，\(X_3\)，\(X_4\)，总分为 \(X\). 对任意一轮，甲乙两人该轮出示的牌等可能地组成 \(16\) 种有序对，其中甲获胜的有 \(6\) 种，所以
\(P(X_k=1)=\frac{6}{16}=\frac38\)（\(k=1\)，\(2\)，\(3\)，\(4\)）.
因此
\(E(X)=\sum_{k=1}^4 E(X_k)=4\times\frac38=\frac32\).
记 \(P(X=j)=p_j\)（\(j=0\)，\(1\)，\(2\)，\(3\)）. 若甲得 \(0\) 分，则只能是甲依次出 \(1\)，\(3\)，\(5\)，\(7\)，乙依次出 \(2\)，\(4\)，\(6\)，\(8\)，故
\(p_0=\frac1{24}\).
若甲得 \(3\) 分，则只能是甲依次出 \(1\)，\(3\)，\(5\)，\(7\)，乙依次出 \(8\)，\(2\)，\(4\)，\(6\)，故
\(p_3=\frac1{24}\).
又因为 \(X\) 的所有可能取值为 \(0\)，\(1\)，\(2\)，\(3\)，所以
\(p_0+p_1+p_2+p_3=1\)，
\(p_1+2p_2+3p_3=E(X)=\frac32\).
代入 \(p_0=\frac1{24}\)，\(p_3=\frac1{24}\) 得
\(p_1+p_2=\frac{11}{12}\)，\(p_1+2p_2=\frac{11}{8}\)，
解得
\(p_2=\frac{11}{24}\).
因此
\[
P(X\ge2)=p_2+p_3=\frac{11}{24}+\frac1{24}=\frac12
\]
故所求概率为 \(\frac12\).
\end{solution}
\section{解答题}

\begin{problem}
记 \(\triangle ABC\) 的内角 \(A\)，\(B\)，\(C\) 对边分别为 \(a\)，\(b\)，\(c\). 已知
\(\sin C =2\cos B\)，\(a^2+b^2-c^2=\sqrt2\,ab\).
\begin{enumerate}
\item 求 \(B\)；
\item 若 \(\triangle ABC\) 的面积为 \(3+\sqrt3\)，求 \(c\).
\end{enumerate}
\end{problem}

\begin{solution}
\begin{enumerate}
\item 由余弦定理
\(a^2+b^2-c^2=2ab\cos C\).
对比可得
\(\cos C=\frac{\sqrt2}{2}\).
又 \(C\in(0,\pi)\)，故
\(\sin C=\sqrt{1-\cos^2C}=\frac{\sqrt2}{2}\).
结合 \(\sin C=2\cos B\)，得 \(\cos B=\frac12\)，故
\(B=\frac\pi3\).

\item 由（1）及 \(\cos C=\frac{\sqrt2}{2}\) 得
\(B=\frac\pi3\)，\(C=\frac\pi4\)，
从而
\(A=\pi-\frac\pi3-\frac\pi4=\frac{5\pi}{12}\).
根据正弦定理，
\(\frac{a}{\sin A}=\frac{b}{\sin B}=\frac{c}{\sin C}\).
故
\(a=\frac{\sin A}{\sin C}c=\frac{\frac{\sqrt6+\sqrt2}{4}}{\frac{\sqrt2}{2}}c=\frac{\sqrt3+1}{2}c\)，\(b=\frac{\sin B}{\sin C}c=\frac{\frac{\sqrt3}{2}}{\frac{\sqrt2}{2}}c=\frac{\sqrt6}{2}c\).
三角形面积
\(S_{\triangle ABC}=\frac12 ab\sin C=\frac{3+\sqrt3}{8}c^2\).
已知面积为 \(3+\sqrt3\)，则
\(\frac{3+\sqrt3}{8}c^2=3+\sqrt3\)，
故
\(c^2=8\)，\(c=2\sqrt2\).
故 \(c=2\sqrt2\).
\end{enumerate}
\end{solution}

\begin{problem}
已知 \(A(0,3)\)，\(P\left(3,\frac32\right)\) 都在椭圆
\(C:\frac{x^2}{a^2}+\frac{y^2}{b^2}=1\ (a>b>0)\)
上.
\begin{enumerate}
\item 求 \(C\) 的离心率；
\item 过 \(P\) 的直线 \(l\) 交 \(C\) 于另一点 \(B\)，且 \(\triangle ABP\) 的面积为 \(9\)，求 \(l\) 的方程.
\end{enumerate}
\end{problem}

\begin{solution}
\begin{enumerate}
\item 由 \(A(0,3)\in C\) 得
\(\frac{3^2}{b^2}=1\)，
所以
\(b^2=9\).
再将 \(P\left(3,\frac32\right)\in C\) 代入得
\(\frac{3^2}{a^2}+\frac{(3/2)^2}{9}=1 \Rightarrow \frac{9}{a^2}+\frac14=1 \Rightarrow a^2=12\).
故离心率
\(e=\sqrt{1-\frac{b^2}{a^2}}=\sqrt{1-\frac{9}{12}}=\frac12\).

\item 直线 \(AP\) 的方程为
\(y=-\frac12x+3\)，
即
\(x+2y-6=0\).
设直线 \(l\) 与椭圆交于另一点 \(B\). 由面积公式
\(S_{\triangle ABP}=9\)
得点 \(B\) 到直线 \(AP\) 的距离为
\(\frac{2S_{\triangle ABP}}{|AP|}=\frac{18}{\frac{3\sqrt5}{2}}=\frac{12}{\sqrt5}\).
设 \(B(x_0,y_0)\). 联立
\(\frac{x_0^2}{12}+\frac{y_0^2}{9}=1\)
及点到直线 \(x+2y-6=0\) 的距离公式
\(\frac{|x_0+2y_0-6|}{\sqrt5}=\frac{12}{\sqrt5}\)，
解得
\(B=(0,-3)\)或\(B=\left(-3,-\frac32\right)\).
于是过 \(P\) 与这两个点的直线分别为
\(3x-2y-6=0\)，\(x-2y=0\).
\end{enumerate}
\end{solution}

\begin{problem}
如图，四棱锥 \(P-ABCD\) 中，\(PA\perp\) 底面 \(ABCD\)，且 \(PA=AC=2\)，\(BC=1\)，\(AB=\sqrt3\).

\begin{enumerate}
\item 若 \(AD\perp PB\)，证明 \(AD\parallel\) 平面 \(PBC\)；
\item 若 \(AD\perp DC\)，且二面角 \(A-CP-D\) 的正弦值为 \(\frac{\sqrt{42}}{7}\)，求 \(AD\).
\end{enumerate}

\bitmapfigure[max width=6.0cm]{img/2024/new_gaokao_paper_1/q17_fig1.png}
\end{problem}

\begin{solution}
\begin{enumerate}
\item 因为 \(PA\perp\) 底面 \(ABCD\)，且 \(AD\subset\) 底面 \(ABCD\)，故 \(PA\perp AD\).
已知 \(AD\perp PB\)，又 \(PA\)，\(PB\) 是平面 \(PAB\) 内两条相交直线，所以 \(AD\perp\) 平面 \(PAB\).
因此 \(AD\perp AB\). 又由
\(AC^2=AB^2+BC^2\)
由勾股定理逆定理得 \(\triangle ABC\) 是以 \(B\) 为直角的直角三角形，从而 \(AB\perp BC\).
因为 \(AD\)，\(BC\) 都在底面 \(ABCD\) 内，且都垂直于 \(AB\)，所以 \(AD\parallel BC\). 又 \(BC\subset\) 平面 \(PBC\)，故 \(AD\parallel\) 平面 \(PBC\).

\item 由 \(AD\perp DC\) 且 \(AC=2\)，设 \(AD=x\)，则
\(CD=\sqrt{4-x^2}\).
过点 \(D\) 作 \(DE\perp AC\) 于 \(E\)，再过点 \(E\) 作 \(EF\perp CP\) 于 \(F\)，连接 \(DF\).
\solutionfigure{\bitmapfigure[max width=3.6cm]{img/2024/new_gaokao_paper_1/q17_solution_fig1.png}}
因为 \(PA\perp\) 底面 \(ABCD\)，所以平面 \(PAC\perp\) 平面 \(ABCD\)，且它们的交线为 \(AC\). 由 \(DE\perp AC\) 且 \(DE\subset\) 平面 \(ABCD\)，得 \(DE\perp\) 平面 \(PAC\).
又 \(EF\perp CP\)，而 \(CP\subset\) 平面 \(PAC\)，所以 \(CP\perp\) 平面 \(DEF\). 因此 \(\angle DFE\) 为二面角 \(A-CP-D\) 的平面角. 已知
\(\sin\angle DFE=\frac{\sqrt{42}}{7}\)，
故
\(\tan\angle DFE=\sqrt6\).
由等面积关系可得
\(DE=\frac{x\sqrt{4-x^2}}{2}\).
又
\[
CE=\sqrt{CD^2-DE^2}=\sqrt{(4-x^2)-\frac{x^2(4-x^2)}{4}}=\frac{4-x^2}{2}
\]
在 \(\triangle PAC\) 中，\(PA=AC=2\)，且 \(PA\perp AC\)，所以 \(\angle ACP=\frac{\pi}{4}\). 又 \(CE\subset AC\)，\(CF\subset CP\)，故 \(\angle ECF=\frac{\pi}{4}\). 因为 \(EF\perp CP\)，所以 \(\angle EFC=\frac{\pi}{2}\)，从而 \(\triangle EFC\) 是等腰直角三角形，所以
\(EF=\frac{CE}{\sqrt2}=\frac{4-x^2}{2\sqrt2}\).
因此
\[
\tan\angle DFE=\frac{DE}{EF}=\frac{\frac{x\sqrt{4-x^2}}{2}}{\frac{4-x^2}{2\sqrt2}}=\frac{x\sqrt2}{\sqrt{4-x^2}}=\sqrt6
\]
解得
\(x^2=3\).
故
\(AD=x=\sqrt3\).
\end{enumerate}
\end{solution}

\begin{problem}
已知函数
\(f(x)=\ln\frac{x}{2-x}+ax+b(x-1)^3\ (x\in(0,2))\).
\begin{enumerate}
\item 若 \(b=0\) 且 \(f'(x)\ge0\)，求 \(a\) 的最小值；
\item 证明：曲线 \(y=f(x)\) 关于点 \((1,a)\) 中心对称；
\item 若 \(f(x)>-2\) 当且仅当 \(1<x<2\)，求 \(b\) 的取值范围.
\end{enumerate}
\end{problem}

\begin{solution}
\begin{enumerate}
\item \(b=0\) 时
\(f'(x)=\frac{2}{x(2-x)}+a\)，\(x\in(0,2)\).
\(f'(x)\) 在 \((0,2)\) 的最小值为 \(2\)（取于 \(x=1\)），故 \(f'(x)\ge0\iff a+2\ge0\)，所以 \(a_{\min}=-2\).

\item 对任意 \(x\in(0,2)\)，有
\(f(2-x)=\ln\frac{2-x}{x}+a(2-x)+b(1-x)^3\)
即
\(f(2-x)=-\ln\frac{x}{2-x}+2a-ax-b(x-1)^3=2a-f(x)\).
因此当点 \((x,f(x))\) 在图像上时，点 \((2-x,2a-f(x))\) 也在图像上，故曲线 \(y=f(x)\) 关于点 \((1,a)\) 中心对称.

\item 若 \(f(1)=-2\)，则 \(a=-2\)（必要条件）. 此时
\(f(x)>-2\iff \ln\frac{x}{2-x}-2(x-1)+b(x-1)^3>0\).
设 \(t=x-1\in(0,1)\)，等价于
\(g(t)=\ln\frac{t+1}{1-t}-2t+bt^3>0\)，\(t\in(0,1)\).
\(g'(t)=\frac{2}{1-t^2}-2+3bt^2=\frac{t^2(-3bt^2+2+3b)}{1-t^2}\).
经讨论可得当 \(b\ge-\frac23\) 时 \(g\) 在 \((0,1)\) 上严格增且 \(g(t)>g(0)=0\)；当 \(b<-\frac23\) 时在 \((0,\frac13)\) 内可取到 \(g'(t)<0\)，不满足全区间正值. 故
\(b\ge-\frac23\).
\end{enumerate}
\end{solution}

\begin{problem}
设 \(m\) 为正整数，数列 \(a_1\)，\(a_2\)，\(\dots\)，\(a_{4m+2}\) 为公差不为 \(0\) 的等差数列，若删去两项 \(a_i\)，\(a_j\ (i<j)\) 后剩余 \(4m\) 项可平均分为 \(m\) 组且每组 4 个数都成等差，则称该数列为 \((i,j)\)-可分数列.
\begin{enumerate}
\item 写出所有 \((i,j)\)，\(1\le i<j\le6\)，使 \(a_1\)，\(\dots\)，\(a_6\) 是 \((i,j)\)-可分数列；
\item 当 \(m\ge3\) 时，证明 \(a_1\)，\(\dots\)，\(a_{4m+2}\) 是 \((2,13)\)-可分数列；
\item 从 \(1\)，\(2\)，\(\dots\)，\(4m+2\) 中不放回任取 \(i<j\)，设 \(a_1\)，\(\dots\)，\(a_{4m+2}\) 是 \((i,j)\)-可分数列的概率为 \(P_m\)，证明 \(P_m>\frac18\).
\end{enumerate}
\end{problem}

\begin{solution}
\begin{enumerate}
\item 可将原等差数列按比例平移缩放为 \(1\)，\(2\)，\(\dots\)，\(4m+2\) 的数列讨论. 对 \(m=1\) 即 \(1\)，\(\dots\)，\(6\)，剩下 4 个数必须是 \(1\)，\(2\)，\(3\)，\(4\) 或 \(2\)，\(3\)，\(4\)，\(5\) 或 \(3\)，\(4\)，\(5\)，\(6\)，故 \((i,j)\) 仅有
\((1,2)\)，\((1,6)\)，\((5,6)\).

\item 对 \(m\ge3\)，去掉 \(2,13\) 后，\(\{1,\dots,4m+2\}\) 中余下 \(4m\) 个数可分成
\(\{1,4,7,10\}\)，\(\{3,6,9,12\}\)，\(\{5,8,11,14\}\)，
及
\(\{15,16,17,18\}\)，\(\{19,20,21,22\}\)，\(\dots\)，\(\{4m-1,4m,4m+1,4m+2\}\)
共 \(m\) 组（当 \(m=3\) 时无第二类）. 故得证.

\item 设
\(A=\{4k+1\mid k=0,1,\dots,m\}\)，\(B=\{4k+2\mid k=0,1,\dots,m\}\).
若 \(i\in A\)，\(j\in B\) 或 \(i\in B\)，\(j\in A\) 且 \(j-i\ne3\)，则该数列一定是 \((i,j)\)-可分. 满足该条件的对数至少有 \((m+1)^2-m\) 个. 总对数 \(\binom{4m+2}{2}=(2m+1)(4m+1)\). 故
\(P_m\ge\frac{(m+1)^2-m}{\binom{4m+2}{2}} =\frac{m^2+m+1}{(2m+1)(4m+1)}>\frac18\).
\end{enumerate}
\end{solution}
